% ===========================================================================
% LDImmersiveEngineering -- the complete manual
%
% Build:  docs/manual/build.sh      (or: python docs/tools/make_manual_data.py
%                                        && latexmk -pdf -cd docs/manual/MANUAL.tex)
% Output: docs/manual/MANUAL.pdf
%
% Two kinds of content live here, and the split is deliberate:
%
%   chapters/   hand-written prose -- what a system is for, how it behaves, why it works the
%               way it does. Written once, revised when the design changes.
%
%   generated/  every table -- blocks, recipes, config, the in-game book -- produced from the
%               source tree by docs/tools/make_manual_data.py. Never edited by hand. A
%               hand-maintained list of blocks is stale the first time somebody adds one, and
%               nobody notices until a player goes looking for something the manual promised.
%
% Regenerating before every build is what "kept in sync with mod updates" actually means here.
% ===========================================================================
\documentclass[11pt,openany]{report}

% ---------------------------------------------------------------------------
% Style for the LDImmersiveEngineering manual.
%
% Everything here is in the MiKTeX base install or its default package set, so a fresh MiKTeX
% will fetch what it needs on the first build and nothing has to be installed by hand.
% ---------------------------------------------------------------------------

\usepackage[T1]{fontenc}
\usepackage[utf8]{inputenc}
\usepackage{lmodern}

\usepackage[a4paper,margin=25mm,headheight=15pt]{geometry}
\usepackage{graphicx}
\usepackage{xcolor}
\usepackage{longtable}
\usepackage{booktabs}
\usepackage{array}
\usepackage{fancyhdr}
\usepackage{titlesec}
\usepackage{enumitem}
\usepackage{tcolorbox}
\usepackage{microtype}
\usepackage[hidelinks]{hyperref}

% ---------------------------------------------------------------------------
% Palette. Taken from the mod's own texture scripts so the book and the blocks agree:
% desaturated gunmetal with one warm accent, which is what IE looks like.
% ---------------------------------------------------------------------------
\definecolor{iesteel}{RGB}{78,78,85}
\definecolor{iesteeldark}{RGB}{40,40,45}
\definecolor{ieorange}{RGB}{233,118,43}
\definecolor{ierule}{RGB}{190,190,196}
\definecolor{iepanel}{RGB}{243,243,244}

% ---------------------------------------------------------------------------
% Headings
% ---------------------------------------------------------------------------
\titleformat{\chapter}[display]
  {\normalfont\sffamily\huge\bfseries\color{iesteeldark}}
  {\color{ieorange}\MakeUppercase{\chaptertitlename}\ \thechapter}{16pt}{\Huge}
\titleformat{\section}
  {\normalfont\sffamily\Large\bfseries\color{iesteeldark}}{\thesection}{1em}{}
\titleformat{\subsection}
  {\normalfont\sffamily\large\bfseries\color{iesteel}}{\thesubsection}{1em}{}

\pagestyle{fancy}
\fancyhf{}
\fancyhead[L]{\sffamily\small\nouppercase{\leftmark}}
\fancyhead[R]{\sffamily\small\thepage}
\renewcommand{\headrulewidth}{0.4pt}

% ---------------------------------------------------------------------------
% Callouts. Three kinds, and the distinction matters:
%   ietip      -- something that makes the mod easier
%   iewarning  -- something that will cost you a build if you get it wrong
%   iefork     -- behaviour that differs from stock Immersive Engineering
% The third is the one this manual exists for; it is coloured to stand out on a flick-through.
% ---------------------------------------------------------------------------
\newtcolorbox{ietip}[1][]{colback=iepanel,colframe=iesteel,fonttitle=\bfseries\sffamily,
  title={Tip},sharp corners,boxrule=0.6pt,#1}
\newtcolorbox{iewarning}[1][]{colback=red!4,colframe=red!55!black,fonttitle=\bfseries\sffamily,
  title={Careful},sharp corners,boxrule=0.6pt,#1}
\newtcolorbox{iefork}[1][]{colback=orange!6,colframe=ieorange,fonttitle=\bfseries\sffamily,
  title={This fork},sharp corners,boxrule=0.6pt,#1}

% ---------------------------------------------------------------------------
% Generated tables.
%
% longtable throughout, because several of these run to many pages and a tabular that cannot
% break would simply be dropped off the bottom of one.
% ---------------------------------------------------------------------------
\newenvironment{blocktable}
  {\small\begin{longtable}{@{}r>{\ttfamily}l p{0.26\textwidth} p{0.40\textwidth}@{}}
   \toprule Meta & \normalfont Registry & Name & Notes \\ \midrule \endhead}
  {\bottomrule\end{longtable}}

\newenvironment{recipetable}
  {\small\begin{longtable}{@{}>{\ttfamily}l r p{0.55\textwidth}@{}}
   \toprule \normalfont Output & Qty & Recipe \\ \midrule \endhead}
  {\bottomrule\end{longtable}}

% One config option: name, type, default, and the comment the mod ships with it.
\newcommand{\configentry}[4]{%
  \noindent\textbf{\ttfamily #1}\hfill{\small\itshape #2, default \texttt{#3}}\\*
  \nopagebreak\small #4\par\vspace{4pt}}

% A block's texture, if the file is there. Missing art must not fail the build -- the manual is
% generated from a tree that is still growing, and a missing PNG is a gap, not an error.
\newcommand{\blockart}[1]{%
  \IfFileExists{../../src/main/resources/assets/immersiveengineering/textures/blocks/#1.png}
    {\includegraphics[width=18mm,height=18mm]{../../src/main/resources/assets/immersiveengineering/textures/blocks/#1.png}}
    {\fbox{\parbox[c][16mm][c]{16mm}{\centering\tiny no art}}}}

% Layer plans for multiblocks.
% ---------------------------------------------------------------------------
% Layer plans for multiblocks.
%
% A multiblock is built one layer at a time, from the bottom up, and that is how these draw it:
% a row of square grids, lowest on the left, each cell carrying a one-letter code for the block
% that goes there. It is how somebody actually builds -- lay a layer, step up, lay the next --
% rather than an isometric picture that looks better and is harder to follow with a stack of
% bricks in your hand.
%
% Drawn rather than screenshotted on purpose. A screenshot of a structure is wrong the moment
% the structure changes and nobody can tell by looking; a plan built from the same letters as
% its own key is at least internally consistent, and is a diff when it changes.
% ---------------------------------------------------------------------------

% ifthen supplies \ifthenelse, which the cell loops below use to test for an empty cell.
\usepackage{ifthen}
\usepackage{tikz}
\usetikzlibrary{positioning,fit,backgrounds}

\definecolor{mbcell}{RGB}{236,236,238}
\definecolor{mbedge}{RGB}{120,120,128}
\definecolor{mbmark}{RGB}{233,118,43}

% \mblayer{caption}{rows}
%
% `rows` is a semicolon-separated list of rows, top row first, each row a comma-separated list
% of cell letters. A dot means "nothing here"; a letter in braces is highlighted, which is how
% the block you strike with the hammer is marked.
%
% Rows are given top-first because that is how you read a plan on paper, and drawn with y
% decreasing so the first row ends up at the top.
\newcommand{\mblayer}[2]{%
  \begin{tikzpicture}[baseline=(current bounding box.north),
      cell/.style={draw=mbedge,fill=mbcell,minimum size=6.5mm,inner sep=0pt,
                   font=\sffamily\scriptsize},
      hot/.style={cell,fill=mbmark!30,draw=mbmark,very thick},
      gap/.style={minimum size=6.5mm,inner sep=0pt}]
    \foreach \row [count=\y] in {#2} {%
      \foreach \c [count=\x] in \row {%
        \ifthenelse{\equal{\c}{.}}
          {\node[gap] at (\x*6.7mm,-\y*6.7mm) {};}
          {\node[cell] at (\x*6.7mm,-\y*6.7mm) {\c};}%
      }%
    }
    \node[font=\sffamily\tiny,anchor=north] at (current bounding box.south) {#1};
  \end{tikzpicture}%
}

% The same, with one cell highlighted: \mblayerhot{caption}{rows}{col}{row}
\newcommand{\mblayerhot}[4]{%
  \begin{tikzpicture}[baseline=(current bounding box.north),
      cell/.style={draw=mbedge,fill=mbcell,minimum size=6.5mm,inner sep=0pt,
                   font=\sffamily\scriptsize},
      hot/.style={draw=mbmark,fill=mbmark!30,very thick,minimum size=6.5mm,inner sep=0pt,
                  font=\sffamily\bfseries\scriptsize},
      gap/.style={minimum size=6.5mm,inner sep=0pt}]
    \foreach \row [count=\y] in {#2} {%
      \foreach \c [count=\x] in \row {%
        \ifthenelse{\equal{\c}{.}}
          {\node[gap] at (\x*6.7mm,-\y*6.7mm) {};}
          {\ifthenelse{\x=#3 \AND \y=#4}
             {\node[hot] at (\x*6.7mm,-\y*6.7mm) {\c};}
             {\node[cell] at (\x*6.7mm,-\y*6.7mm) {\c};}}%
      }%
    }
    \node[font=\sffamily\tiny,anchor=north] at (current bounding box.south) {#1};
  \end{tikzpicture}%
}

% A row of layers with a shared key underneath.
%
% \begin{mbplan}{Coke Oven}{B = Coke Brick}
%   \mblayer{Layer 1}{B,B,B;B,B,B;B,B,B} \hfill ...
% \end{mbplan}
\newenvironment{mbplan}[2]
  {\par\vspace{6pt}\noindent\textbf{\sffamily\small #1}\\[2pt]%
   \def\mbplankey{#2}\noindent}
  {\\[2pt]{\footnotesize\itshape\mbplankey}\par\vspace{8pt}}


\setlist{noitemsep,topsep=4pt}
\graphicspath{{figures/}}


\title{\Huge\sffamily\bfseries LDImmersiveEngineering\\[6pt]
       \Large\mdseries The Complete Manual}
\author{\sffamily A private fork of Immersive Engineering for Minecraft 1.12.2}
\date{\sffamily\today}

\begin{document}
\maketitle

\begin{abstract}
\noindent
This is the whole of Immersive Engineering as this fork ships it: every block, every machine,
every system, and the four large additions that are not in the mod upstream --- the virtual power
grid, the petroleum chain, the virtual fluid network, and conduits --- together with the city-mode
performance work that runs underneath all of it.

Where this fork differs from stock Immersive Engineering, the difference is called out in an
orange box. Where a number appears, it was read out of the source rather than remembered.
\end{abstract}

\tableofcontents

% ---------------------------------------------------------------------------
% Part I -- getting started, and the mod's own idea of itself
% ---------------------------------------------------------------------------
\part{The mod}
\chapter{What this is}
\label{ch:introduction}

Immersive Engineering is a technology mod built around the idea that machinery should look like
machinery. Its power is carried on wires strung between poles, its ore is crushed by a drum you can
watch turn, and the large machines are structures you assemble block by block rather than single
cubes that happen to contain a factory.

This manual documents \textbf{LDImmersiveEngineering}, a private fork of that mod for Minecraft
1.12.2. Everything the original does, this does. On top of it sit four large additions and one
performance regime, none of which exist upstream:

\begin{description}
\item[The virtual power grid] Moves Immersive Flux between places with no wire between them.
      \autoref{ch:grid}.
\item[Petroleum] A complete oil chain: prospect, drill, refine, store, burn, sell.
      \autoref{ch:petroleum}.
\item[The virtual fluid network] The same idea as the grid, for fluid. \autoref{ch:fluidnet}.
\item[Conduits] Surface-mounted bundled wiring for the inside of buildings. \autoref{ch:conduits}.
\item[City mode] Trades simulation detail for server tick time, subsystem by subsystem.
      \autoref{ch:citymode}.
\end{description}

\section{How to read this}

Part~II is Immersive Engineering as it ships. Part~III is what the fork adds. Part~IV is generated
straight from the source tree --- every block, every recipe, every configuration option --- and is
the place to look when you want a number rather than an explanation.

Three kinds of callout appear throughout:

\begin{ietip}
Something that makes the mod easier, or a habit worth forming early.
\end{ietip}

\begin{iewarning}
Something that will cost you a build, a machine or an afternoon if you get it wrong.
\end{iewarning}

\begin{iefork}
Behaviour that differs from stock Immersive Engineering. If you already know the mod, these boxes
are the parts you have to read.
\end{iefork}

\section{Where the numbers come from}

Every figure in this manual was read out of the source. Where a number is configurable --- and most
are --- the value quoted is the shipped default, and \autoref{ch:config} lists the option that
changes it.

\begin{iewarning}
A modpack may have changed any of them. If a machine in your world does not match a number here,
check the config before assuming the manual is wrong.
\end{iewarning}

\chapter{Getting started}
\label{ch:getting-started}

Immersive Engineering does not begin with a machine. It begins with a hammer and a fire.

\section{The Engineer's Hammer}

Three iron ingots and two sticks. It is the mod's universal tool and you will never stop carrying
one:

\begin{itemize}
\item \textbf{Assembles multiblocks.} Every structure in \autoref{ch:multiblocks} is built by
      placing its blocks and then striking one of them.
\item \textbf{Rotates blocks.} Right-click most machines to turn them.
\item \textbf{Configures sides.} Sneak-right-click a machine face to cycle what that side does.
\item \textbf{Dismantles.} Takes a multiblock apart into its parts.
\end{itemize}

\section{The Engineer's Manual}

One book and one lever. It is the in-game version of this document, and it updates itself as you
obtain things. \autoref{ch:ingame} reproduces it in full.

\section{Treated wood}

Creosote oil in a barrel or a Bottling Machine turns planks into \textbf{treated wood}, which
resists weather and is the base of almost everything early. You get creosote from the coke oven, so
the coke oven comes first.

\section{The Coke Oven}

A $3\times3\times3$ of Coke Bricks. Turns coal into \textbf{coal coke} and \textbf{creosote oil},
both of which you need continuously:

\begin{itemize}
\item Coal coke burns longer than coal and is the fuel for the Blast Furnace.
\item Creosote oil treats wood and, later, is a Squeezer input.
\end{itemize}

\begin{ietip}
Build two. The oven is slow, both outputs are consumed constantly, and the second one costs almost
nothing next to the time it saves.
\end{ietip}

\section{The Blast Furnace}

A $3\times3\times3$ of Blast Bricks. Iron plus coal coke gives \textbf{steel}, which is the material
the entire middle of the mod is built from. It also produces \textbf{slag}, a Blast Brick
ingredient, so the furnace partly pays for its own expansion.

The \textbf{Advanced Blast Furnace} is the upgrade: faster, and it accepts external heating from a
Blast Furnace Preheater.

\section{Your first power}

Immersive Engineering measures power in \textbf{Immersive Flux (IF)}, which is Forge Energy under
another name --- it interoperates with RF and FE at one to one.

The cheapest start is a \textbf{Water Wheel} or a \textbf{Windmill}: both are built from treated
wood, need no fuel and produce a trickle. Put a \textbf{Dynamo} behind the wheel, an \textbf{LV
Connector} on the dynamo, and run copper wire to whatever needs it.

\begin{iewarning}
Wires are not decorative. Walking into an energised HV wire will hurt you, and connecting a wire to
a connector of the wrong tier will destroy it. See \autoref{ch:wiring}.
\end{iewarning}

\section{The order most people build in}

\begin{enumerate}
\item Hammer, manual, coke oven.
\item Treated wood, then a water wheel or windmill for a first trickle of power.
\item Blast furnace, then steel.
\item Crusher, for doubled ore.
\item Metal Press, for cheap plates and rods --- this is the point the mod starts paying you back.
\item Diesel generator or one of the fork's power plants (\autoref{ch:petroleum}) for real output.
\end{enumerate}


% ---------------------------------------------------------------------------
% Part II -- stock Immersive Engineering
% ---------------------------------------------------------------------------
\part{Immersive Engineering}
\chapter{Materials}
\label{ch:materials}

\section{Ores}

Immersive Engineering adds \textbf{copper}, \textbf{aluminium} (bauxite), \textbf{lead},
\textbf{silver}, \textbf{nickel} and \textbf{uranium}. All generate in ordinary stone; uranium is
deep and sparse.

Every one of them can be doubled in a Crusher and, with the Arc Furnace, taken further still.

\section{Alloys}

Made in the Alloy Smelter, or in an Arc Furnace once you have one:

\begin{description}
\item[Constantan] Copper + nickel. Used in the Thermoelectric Generator and MV components.
\item[Electrum] Silver + gold. The MV wire metal.
\item[Steel] Iron + coal coke, in the Blast Furnace rather than the Alloy Smelter. The HV wire
      metal and the backbone of everything.
\end{description}

\section{Plates, rods and wire}

Three shapes, and the machines that make them cheaply are the reason to build those machines:

\begin{description}
\item[Plates] Hammered from ingots by hand, or pressed far more cheaply in a \textbf{Metal Press}
      with a plate mould.
\item[Rods (sticks)] Same, with a rod mould.
\item[Wire] Same, with a wire mould. Wire is what wire coils are made of.
\end{description}

\begin{ietip}
Hand-hammering plates works and is miserable at volume. The Metal Press pays for itself within one
build, and its moulds are the cheapest thing in the mod to duplicate.
\end{ietip}

\section{Sheet metal}

Decorative and structural blocks in every metal the mod knows. They are the visual vocabulary of an
IE build --- most large multiblocks are faced in them --- and several structures require them
outright.

\section{Concrete}

Sand, gravel, clay and water give \textbf{concrete}, which sets from a liquid. It is blast
resistant, comes in several finishes, and is what most players end up building the actual buildings
from.

\section{Hemp}

A crop grown from Hemp Seeds. Two blocks tall, harvested from the top. Gives \textbf{plant fibre},
which becomes \textbf{fabric} and then the cloth used in Balloons, Shaders and the Powerpack.

\section{Coal coke and creosote}

Covered in \autoref{ch:getting-started}, and worth repeating because they never stop being needed:
coke fuels the Blast Furnace and creosote treats wood.

\section{Fluids}

Stock Immersive Engineering ships creosote, plant oil, ethanol, biodiesel and concrete.

\begin{iefork}
This fork adds the entire petroleum family --- crude, the refined fractions, and the products of
cracking. \autoref{ch:petroleum}.
\end{iefork}

\chapter{Wiring}
\label{ch:wiring}

Power in Immersive Engineering travels along wires strung between connectors. The wires sag, they
have length limits, they lose energy over distance, and at high voltage they will kill you. That is
the whole character of the mod in one system.

\section{The three tiers}

\begin{center}
\small
\begin{tabular}{@{}llrrr@{}}
\toprule
Tier & Wire & Max transfer & Loss & Max length \\
\midrule
LV & Copper    & 2\,048 IF/t  & 5\,\%   & 16 blocks \\
MV & Electrum  & 8\,192 IF/t  & 2.5\,\% & 16 blocks \\
HV & Steel     & 32\,768 IF/t & 2.5\,\% & 32 blocks \\
\bottomrule
\end{tabular}
\end{center}

Loss is proportional to how much of the wire's maximum length the run actually uses, so a short HV
run loses very little and a run at the limit loses the full percentage.

\begin{ietip}
Insulated copper and insulated electrum carry the same power as their bare versions and do not shock
anything that touches them. They cost more. On a build people walk through, pay it.
\end{ietip}

\section{Connectors and relays}

A \textbf{connector} is where a wire meets a machine: it takes power from the block it is attached
to, or gives power to it. A \textbf{relay} only passes wire to wire --- it is a pole, not a socket
--- and relays are how you get a line across a distance without a machine at every corner.

\begin{iewarning}
A connector accepts only its own tier's wire. Attaching the wrong one destroys the connector.
\end{iewarning}

\section{Transformers}

Step between tiers. An \textbf{LV/MV Transformer} joins a low-voltage side to a medium-voltage one;
an \textbf{MV/HV Transformer} does the same higher up. Wire the correct tier to each side --- they
are not symmetrical, and the block's model tells you which side is which.

\section{Wire coils and the wirecutter}

Wire is placed by holding a \textbf{coil} and right-clicking one connector, then the other. The
\textbf{Engineer's Wirecutter} removes a wire and returns the coil.

\section{Posts and arms}

\textbf{Wooden} and \textbf{Steel Posts} are what you actually string wire between across country.
Both accept arms, which give you somewhere to mount connectors and transformers off the pole itself.

\section{Breaker switches and meters}

\begin{description}
\item[Breaker Switch] A manual on/off in a wire run, also switchable by redstone.
\item[Redstone Breaker] The same, driven by redstone alone.
\item[Energy Meter] Reports throughput on the line passing through it.
\item[Current Transformer] Emits a redstone signal proportional to the power flowing.
\end{description}

\section{Feedthroughs}

An insulator that carries a wire through a wall without breaking the run, so a line can enter a
building without a connector on each side.

\section{Wire damage}

An energised HV wire damages anything that touches it, and an overloaded wire burns out and drops
as scrap.

\begin{iefork}
City mode changes both. Wire burnout is disabled outright, and shock damage becomes local: rather
than every powered connector broadcasting to the whole network every tick on the chance somebody is
touching a wire, a node that needs the information asks for it. That single broadcast was the most
expensive thing the mod did on a profiled server. \autoref{ch:citymode}.
\end{iefork}

\begin{iefork}
This fork also offers two alternatives to wire entirely: the virtual grid (\autoref{ch:grid}) for
distance, and conduits (\autoref{ch:conduits}) for indoors.
\end{iefork}

\chapter{Generators}
\label{ch:generators}

\section{Water Wheel}

Treated wood. Mounted on a \textbf{Dynamo}, turned by flowing water down one side. Output scales
with how much of the wheel is being driven, and three wheels can share one dynamo.

Cheap, silent, needs no fuel, produces very little. It is the right first generator and never the
last.

\section{Windmill}

Treated wood on a dynamo again, mounted high and in the open. Output rises with altitude and falls
if obstructed. The \textbf{Improved Windmill} adds iron sheet metal for a substantial increase.

\begin{ietip}
Windmills need clear air around them, not merely above them. A windmill in a valley is a decoration.
\end{ietip}

\section{Thermoelectric Generator}

Two blocks of differing temperature on opposite sides --- classically lava and packed ice, or a
blaze-adjacent block against a cold one --- and it produces a steady trickle with no moving parts
and no fuel logistics. The default output multiplier is~1.

Its virtue is that it never stops. Its vice is that it is slow.

\section{Diesel Generator}

The first serious generator: a multiblock burning biodiesel for \textbf{4\,096 IF/t} by default.
Needs a fluid supply, and rewards a real one.

\begin{iefork}
The petroleum chain supplies it far better than the biodiesel loop does, and adds several larger
plants above it, ending at the Gas Turbine. \autoref{ch:petroleum}.
\end{iefork}

\section{Lightning Rod}

A tall multiblock that captures strikes for \textbf{16\,000\,000 IF} apiece by default. Entirely
weather-dependent, spectacular, and useless as a base load --- pair it with storage large enough to
hold a strike or most of it is wasted.

\section{Storage}

\begin{description}
\item[LV / MV / HV Capacitor] Buffers of increasing size. Each face can be configured to input,
      output or nothing with a hammer.
\item[Capacitor Backup] Keeps a capacitor's contents through a chunk unload.
\end{description}

\begin{iewarning}
A capacitor with every face set to output will happily drain itself into whatever is adjacent.
Configure the faces deliberately.
\end{iewarning}

\begin{iefork}
City mode makes fuel cosmetic for generators: a presence check and a token sip instead of a per-tick
burn rate and tank drain. A generator with fuel in it runs; how much is left stops being tracked.
\end{iefork}

\chapter{Machines}
\label{ch:machines}

The single-block devices. Everything here is placed, powered and used on its own; the assembled
structures are in \autoref{ch:multiblocks}.

\section{Powered}

\begin{description}
\item[Sample Drill] Takes a core sample of the chunk, telling you what mineral vein is beneath it.
      The prerequisite for an Excavator.
\item[Charging Station] Charges tools, armour and any item that holds flux.
\item[Fluid Pump] Moves fluid through pipes, and is what makes a pipe network do anything.
\item[Fluid Placer] Places fluid into the world as a source block.
\item[Garden Cloche (Belljar)] Grows a crop in a glass jar from water and fertiliser. Slow, tidy,
      and entirely automatic.
\item[Floodlight] A directional lamp. Expensive to run and the brightest light in the mod.
\item[Turret (Gun / Chemical)] An automated defence, configurable by owner and target type.
\item[Tesla Coil] Damages entities in a radius. Genuinely dangerous to its owner.
\item[Electric Lantern] A small lamp. Cheap enough to use in quantity.
\end{description}

\begin{iefork}
City mode caps how many light blocks a Floodlight may place and stops it re-tracing its beams on a
timer --- it re-traces only when the light switches or a neighbour changes. Floodlights are usually
the most expensive block in a city build.
\end{iefork}

\section{Unpowered}

\begin{description}
\item[Wooden Barrel] Holds fluid, and is where creosote treats wood by hand.
\item[Wooden Crate] Storage that keeps its contents when broken.
\item[Reinforced Crate] The same, lockable.
\item[Workbench] Where blueprints are used and tools are modified.
\item[Item Router] Sorts items by filter into up to six outputs.
\item[Fluid Sorter] The same for fluid.
\item[Redstone Wire Connector] Carries a redstone signal down a wire rather than along the ground.
\end{description}

\section{Fluid pipes}

Steel pipes carrying fluid between machines. They connect automatically, can be covered with sheet
metal to hide them, and need a Fluid Pump to move anything.

\begin{iefork}
Two changes here. The pipe network's flood fill was rewritten --- it used a linear scan for its
closed set and removed from the head of a list for its open set, which is quadratic on a large tank
farm --- and city mode lets a pipe hand its fluid to endpoints in order rather than simulating a
fill against every one of them and splitting in proportion. Nothing is created or destroyed; only
which of several tanks fills first changes.
\end{iefork}

\chapter{Multiblocks}
\label{ch:multiblocks}

Every structure here is built by placing its blocks in the right shape and then striking it with an
Engineer's Hammer. The in-game manual renders each one as a rotating diagram; \autoref{ch:ingame}
carries the text, and the exact block lists are in \autoref{ch:blocks}.

\begin{ietip}
Build multiblocks with the hammer already in your hand and the manual page open. Almost every failed
assembly is one block in the wrong place, and the diagram will show you which.
\end{ietip}

\section{Processing}

\begin{description}
\item[Coke Oven] $3\times3\times3$. Coal to coal coke and creosote oil. Needs no power.
\item[Blast Furnace] $3\times3\times3$. Iron to steel, with slag. Needs no power.
\item[Advanced Blast Furnace] Faster, and accepts a Preheater for more speed still.
\item[Alloy Smelter] Combines two metals. Needs no power.
\item[Crusher] Doubles ore, and grinds a great many other things. The first machine most players
      build that visibly pays for itself.
\item[Arc Furnace] Recycles, alloys and processes ores further than the Crusher can. Consumes
      graphite electrodes, which wear out.
\item[Metal Press] Ingots to plates, rods and wire with a mould. Cheap to run and enormously
      cheaper than hammering.
\item[Squeezer] Presses seeds to plant oil, and other things besides.
\item[Fermenter] Ferments crops to ethanol.
\item[Refinery] Plant oil and ethanol to biodiesel.
\item[Mixer] Combines fluids with items --- concrete, most importantly.
\item[Bottling Machine] Fills containers with a fluid.
\item[Auto Workbench] Automates blueprint crafting.
\item[Assembler] Automates ordinary crafting from up to three recipes at once, pulling from a
      shared inventory.
\end{description}

\section{Mining}

\begin{description}
\item[Excavator] The largest structure in the mod. Digs a mineral vein located by Sample Drill,
      producing ore in quantity until the vein depletes. Consumes \textbf{4\,096 IF/t} by default.
\item[Bucket Wheel] The visible half of the Excavator, and the reason it is worth building where
      people can see it.
\end{description}

\section{Storage and power}

\begin{description}
\item[Sheetmetal Tank] A $5\times5\times5$ vessel holding 512\,000\,mB.
\item[Silo] The same idea for solid items.
\item[Diesel Generator] \autoref{ch:generators}.
\item[Lightning Rod] \autoref{ch:generators}.
\end{description}

\begin{iefork}
The fork adds several: the Grid Management Console and Substation (\autoref{ch:grid}), the Fluid
Control Console (\autoref{ch:fluidnet}), and the whole petroleum chain --- Derrick, Cracking Unit,
Loading Gantry, the buried tanks and the free-form Storage Tank (\autoref{ch:petroleum}).
\end{iefork}

\begin{iefork}
City mode stops idle multiblocks re-scanning the recipe list every tick and widens the scan interval
for the rest. A room full of idle machines is close to free.
\end{iefork}

\chapter{Logistics}
\label{ch:logistics}

\section{Conveyor belts}

The mod's item transport, and one of the things it is known for: items ride visibly on top rather
than vanishing into a pipe.

\begin{description}
\item[Conveyor Belt] Moves items along. Right-click with a hammer to change direction.
\item[Uncontrolled] Does not stop items at the end --- they fall off.
\item[Dropper] Drops items through the block beneath.
\item[Vertical] Carries items upward.
\item[Extract / Insert] Pull from and push into adjacent inventories.
\item[Splitter] Alternates output between two directions.
\item[Chute] A vertical shaft. Fast, and made of any sheet metal.
\end{description}

\begin{ietip}
Belts can be covered with any block, so a line can run through a wall or under a floor and still be
watched where you want it seen.
\end{ietip}

\section{Sorting}

\begin{description}
\item[Item Router] Six outputs, each with its own filter. The workhorse.
\item[Fluid Sorter] The same, for fluid.
\end{description}

\section{Storage}

\begin{description}
\item[Wooden Crate] Keeps its contents when broken, so it doubles as a moving box.
\item[Reinforced Crate] Lockable to its owner.
\item[Toolbox] A portable container for tools, food and ammunition.
\end{description}

\section{Posts and scaffolding}

\textbf{Wooden} and \textbf{Steel Scaffolding} climb like ladders and are the mod's structural
vocabulary. \textbf{Posts} carry wire across country and take arms for mounting hardware.

\section{Redstone}

\begin{description}
\item[Redstone Wire Connector] Carries a redstone signal along a wire.
\item[Breaker Switch] A lever in a power line.
\item[Current Transformer] A redstone signal proportional to power flowing.
\end{description}

\begin{iefork}
Conduits carry redstone too, on any of their sixteen channels, alongside power in the same bundle.
\autoref{ch:conduits}.
\end{iefork}

\chapter{Tools and equipment}
\label{ch:tools}

\section{The essentials}

\begin{description}
\item[Engineer's Hammer] Assembles, rotates, configures and dismantles. Never stop carrying one.
\item[Engineer's Wirecutter] Removes wire and returns the coil.
\item[Engineer's Voltmeter] Reads a machine's stored energy and throughput.
\item[Engineer's Manual] The in-game book.
\end{description}

\begin{iefork}
The Voltmeter gained a second job: sneak-click a grid box to pick up its segment, sneak-click others
to assign them, sneak-click the air to clear. There is no Engineer's Screwdriver in 1.12 --- it is a
later item --- so the quick-assign gesture lives here instead.
\end{iefork}

\section{Powered tools}

\begin{description}
\item[Engineer's Drill] A mining tool taking a drill head, which wears out and can be replaced.
      Mines in a $3\times3$ with the right upgrade and is the fastest miner in the mod.
\item[Buzzsaw] Fells trees, and cuts wood into planks as it goes.
\item[Chemical Thrower] Sprays a fluid --- extinguishing, igniting or poisoning depending on what
      you load.
\item[Railgun] Fires rods at very high velocity for very high damage.
\item[Revolver] The mod's sidearm, with a wide range of ammunition and cosmetic parts. Extensively
      upgradeable at a Workbench.
\end{description}

\section{Upgrades}

Most powered tools take upgrades at a Workbench: capacity, speed, extra barrels, larger tanks,
scopes. The upgrade is an item and can be removed again.

\section{Armour and equipment}

\begin{description}
\item[Steel Armour] A full set, better than iron.
\item[Faraday Suit] Protects against electricity --- including the mod's own wires and Tesla Coils.
\item[Powerpack] A back-worn battery charging held tools as you use them.
\item[Earmuffs] Silence machinery near you.
\item[Shield] Blocks, and takes upgrades of its own.
\item[Climbing Rope, Skyhook] Movement. The Skyhook rides along wire runs, which is the most
      enjoyable transport in the mod once you have a wire network worth riding.
\end{description}

\section{Shaders}

Purely cosmetic skins for tools, armour and some machines, obtained from Shader Bags. They do not
change behaviour and there are a great many of them.

\begin{iefork}
One held item is new: the Engineer's Network Terminal, which reads the power grid and the fluid
network from anywhere and changes neither. \autoref{ch:grid}.
\end{iefork}


% ---------------------------------------------------------------------------
% Part III -- what this fork adds
% ---------------------------------------------------------------------------
\part{This fork}
\chapter{The Virtual Power Grid}
\label{ch:grid}

Running catenary wire across a city is honest work and is not always practical. The
\textbf{virtual power grid} moves Immersive Flux between places that are not connected by any wire
at all --- across a mountain, across a chunk border, across dimensions.

\begin{iefork}
Not in stock Immersive Engineering. It takes the idea Flux Networks made popular and rebuilds it in
IE's idiom: sheet-metal service boxes, a console you walk to, and named segments rather than
per-player networks.
\end{iefork}

\section{Feed Units and Service Units}

Power still has to enter and leave somewhere. It enters at a \textbf{Grid Feed Unit} and leaves at a
\textbf{Grid Service Unit}. Both bolt to any solid surface, including engineering posts.

These boxes exchange power with the block they are bolted to, exactly as a machine does. They do
\emph{not} accept wires directly --- put a connector against one if you want it to feed a wire
network. A Feed Unit takes power out of the world; a Service Unit puts it back. Both carry a
terminal post on the front, which is where wiring goes.

\begin{ietip}
A wire connector \emph{beside} a Service Unit is fed as well as one bolted to it, whichever way the
connector faces. IE connectors normally take power on one side only --- the block they are mounted
on --- and a connector on the wall next to a unit would otherwise sit there touching live hardware
and doing nothing.
\end{ietip}

\section{Segments}

The grid is not one pool. It is divided into \textbf{segments}: named groups of boxes, each with its
own switch, its own transfer limits and its own statistics. A segment is the thing you switch off
when you want the east side of town dark; the boxes on it are the sockets.

A box that has not been assigned to a segment sits with an unlit lamp doing nothing. Assigning it is
the whole configuration step.

Each segment carries a \textbf{priority} order and a \textbf{transfer cap} per device, so you decide
what is served first when there is not enough to go round. Anything marked \textbf{critical} is
served before everything else --- the pumps, the blast furnace, the hospital.

\section{The Grid Management Console}

A $2\times2$ wall of four different blocks, struck anywhere on its front face with an Engineer's
Hammer:

\begin{center}
\small
\begin{tabular}{@{}lll@{}}
\toprule
 & Left & Right \\
\midrule
Upper & Console Housing (the terminal) & Redstone Engineering Block \\
Lower & Light Engineering Block & Heavy Engineering Block \\
\bottomrule
\end{tabular}
\end{center}

\noindent The terminal is the screen, the redstone block is the instrument rack beside it, and the
light and heavy blocks are the desk and the power cabinet under them. The finished console is drawn
as one piece --- one desk, one monitor spanning both blocks --- and breaking it returns the four
components.

Six tabs: an overview of every segment, the segment editor, the device list, the failover editor,
live statistics, and settings. Right-clicking any box also opens its own small panel, so a freshly
placed unit can be assigned without walking back.

\section{Failover}

A segment may name another as its \textbf{backup}. If the first cannot meet its load, the second
covers it. With \emph{top-up} enabled a backup also covers ordinary shortfalls, not only outages;
with it off, backups engage on outage or trip alone.

\section{Breakers and schedules}

If breakers are enabled, a segment held at its output ceiling long enough will \textbf{trip} and
latch off until somebody flips it back on. That is a real switch on the Segments tab.

A segment may also carry a \textbf{schedule}: two times of day between which it may run. The
schedule is a gate, never a second switch --- it can hold a segment down but never raise one a
player switched off, because otherwise the console toggle and the clock fight every dusk and
whichever ran last wins.

\section{The Grid Signal Unit}

Bridges a segment to redstone in either direction. In \emph{input} mode a signal holds its segment
off, which makes any lever, daylight sensor or logic circuit an external kill switch; inverted, it
demands a keep-alive signal instead, which is a dead-man's switch. In \emph{output} mode it reports
whether its segment is up --- or, inverted, whether it is down, which is a fault lamp.

\section{The Engineer's Network Terminal}

A held item showing every segment and every fluid main, from anywhere, and changing neither.

\begin{ietip}
The terminal is read-only by construction rather than by convention --- there is no path from it to
a change, whatever the client sends. Changes are made at a console.
\end{ietip}

\section{Chunk loading}

A device may be marked to keep its chunk loaded, within a server-wide budget. Fittings sharing a
chunk cost one between them. A disabled device stops holding its chunk.

\section{City mode}

See \autoref{ch:citymode}. In short: segments stop accounting for flux and switch to presence. A
segment is energised or it is not, and its Service Units deliver freely.

\chapter{Petroleum}
\label{ch:petroleum}

A complete oil chain: find it, drill it, refine it, store it, burn it, and sell it at a forecourt.

\begin{iefork}
Entirely this fork, and the largest single addition in it.
\end{iefork}

\section{Prospecting}

Oil sits in \textbf{reservoirs} beneath the world, which have to be found before they can be
drilled. The survey kit reports bearings and ranges rather than handing you a coordinate.

\section{Drilling}

A \textbf{Derrick} sunk over a reservoir brings crude to the surface through a \textbf{Wellhead}. A
reservoir is finite: its flow decays as it empties, and it can be blown out by mishandling.

\section{Refining}

Crude is separated into fractions, and a \textbf{Cracking Unit} converts between them. The cracker
has no temperature dial --- the dial is the firebox. Run it hot on heavy fuel oil and it yields
gasoline; run it cool on gas and it yields diesel. A knob built out of blocks.

\section{Storage}

Three \textbf{buried tanks} --- Domestic, Commercial and Bulk --- are invisible except for a surface
fill cap carrying the gauge. Build the shell, cover the roof, and hammer the cap.

The \textbf{Storage Tank} is the other kind: a hollow box of any rectangular shape from
$3\times3\times3$ to $16\times16\times16$, holding sixteen buckets per enclosed block. Build the
shell you want and strike any wall with a hammer. Pipe into any wall block, not just one fitting.

\begin{ietip}
Buried tanks are for supply you do not want to look at --- under a forecourt, under a yard. The
Storage Tank is for a tank farm you do.
\end{ietip}

For propane specifically there are three single-block vessels: the \textbf{Propane Cylinder} at four
buckets, the \textbf{Upright Propane Tank} at twelve, and the \textbf{Torpedo Propane Tank} at
twenty-four. They differ only in size and shape --- all three take propane and nothing else, all
three keep their contents when broken, and all three fill from a bucket or a pipe. Pick the one that
suits what you are building: a barbecue, a workshop wall, or the end of a garden.

\section{Burning it}

A ladder of power plants ending at the Gas Turbine, plus the \textbf{Portable Generator} for a build
site.

\section{The forecourt}

A \textbf{Gas Station Pump}, a \textbf{Forecourt Price Sign} and a nozzle that binds to a pump
within eight blocks. Fuel is dispensed to a person rather than to a machine.

\begin{iewarning}
Nothing past the Gas Turbine has been extensively playtested. The arithmetic is unit-tested and the
server starts clean, but the balance of the later chain is still unproven.
\end{iewarning}

\chapter{The Virtual Fluid Network}
\label{ch:fluidnet}

What the virtual grid is for power, this is for fluid: it moves a fluid between places with no pipe
between them.

\begin{iefork}
Not in stock Immersive Engineering. It is deliberately built as a mirror of the power grid rather
than as a generalisation of it --- the two have separate save data, so a mistake in one cannot cost
you the other.
\end{iefork}

\section{Mains}

The fluid network is divided into \textbf{mains}, which are what segments are to the grid: named
groups of fittings with their own switch, limits and statistics.

A main carries \textbf{one fluid}. The first inlet to deliver decides which; after that the main
refuses anything else until it drains. That is not a limitation to work around --- it is what makes
a main something you can name ``diesel'' and trust.

\section{Fittings}

\begin{description}
\item[Fluid Inlet] Takes fluid from an adjacent tank or pipe and puts it into its main.
\item[Fluid Outlet] Draws from its main and fills whatever sits against it. One marked
      \textbf{critical} keeps its supply when the main runs short.
\item[Main Valve] Shuts a main by hand or on a redstone signal, and reports whether it is
      pressurised.
\item[Fluid Control Console] Where the network is managed. A separate structure from the power
      console, because the two networks fail apart.
\end{description}

\section{Failover and schedules}

As the grid: a main may name a backup, and may carry a schedule that can hold it closed but never
open one somebody shut.

\section{City mode}

Mains stop accounting for millibuckets and switch to presence --- a main is pressurised or it is
not.

\chapter{Conduits}
\label{ch:conduits}

Wire sags. Strung across a valley that is exactly right; strung along a ceiling it is exactly wrong.
\textbf{Conduit} is the indoor answer: thin tubing that clips flat against whatever face you place
it on and runs in that face, turning in right angles.

\begin{iefork}
Not in stock Immersive Engineering. It takes the bundled-cable idea from SimpleLogic and puts it in
IE's idiom --- surface conduit, junction boxes, and ordinary IE connectors at the breakouts.
\end{iefork}

One length of conduit carries \textbf{sixteen} independent conductors, named after the sixteen dyes.

\section{Runs}

Place conduit against a wall, floor or ceiling and it clips to whatever you clicked. Lengths joined
along the same surface make a \textbf{run}.

To continue a run, click the end of it. A conduit placed against another conduit joins that
conduit's surface rather than clipping to the conduit itself, so you can lay a corridor by clicking
along the run instead of hunting for the bare wall between lengths. The first length off a junction
box does the same thing --- it looks for the surface behind it rather than clipping to the box.

\textbf{A run stays on one surface.} Conduit on a floor does not climb the wall by itself. That is
deliberate: a plane change goes through a junction box, which is a real block and what you would put
at a corner anyway.

Runs are discovered, not drawn. Laying conduit between two boxes creates the connection; pulling a
length out of the middle breaks it. There is no coil and no linking tool --- the run is the thing
you built.

\begin{mbplan}{A run seen from above, on a floor}{C = conduit, J = junction box. The run turns in
the plane of the surface it is clipped to; leaving that surface needs the box.}
\mblayer{}{J,C,C,C,C,J;.,.,.,.,.,C;.,.,.,.,.,J}
\end{mbplan}

\section{The Junction Box}

A patch panel with six faces. Each face is one of three things:

\begin{itemize}
\item Right-click with a \textbf{dye}: that face breaks out the conductor of that colour.
\item Right-click a patched face with \textbf{redstone dust}: cycles it between power, reading
      redstone, and emitting redstone.
\item Sneak and right-click empty-handed: unpatches the face.
\end{itemize}

\textbf{A patched face wears a plate in its conductor's colour}, so a box tells you how it is wired
from across the room: which faces break out, and which channel each one carries. A box with no
plates on it is a box nobody has patched.

Put an LV, MV or HV Connector against a power breakout and it picks up that conductor and nothing
else. One bundle down a corridor can feed a lighting circuit, a workshop and a furnace, each on its
own colour.

\begin{ietip}
There is no tier setting. The tier of a circuit is whichever connector you hang on its breakout,
because IE's connectors already cap throughput by tier. Mixed tiers in one bundle are fine.
\end{ietip}

A colour lives on one face at a time --- patching it somewhere new takes it off wherever it was,
because the same conductor arriving at two connectors is a short.

A box passes the whole bundle through whatever is patched, so one dropped in to turn a corner or
change surface needs no configuration at all.

\section{The Ground Feeder}

A junction box is the right way to leave a surface indoors, where it reads as hardware bolted to a
wall. Out in a field it is a steel box sitting in a lawn, and there is no surface to change
\emph{to} until the run is already underground.

The \textbf{Ground Feeder} is the other answer: a whole cube that fills a hole in a floor, a wall or
the ground, passes a bundle straight through along one axis, and \textbf{draws itself as whatever is
around it}.

\textbf{It looks at what is visible rather than at what is merely nearby.} A feeder set into a grass
field has eight grass blocks around it and nine dirt blocks underneath, so counting neighbours would
put dirt in the middle of a lawn. Instead a block scores triple if any face of it can be seen at all,
and double if it shares a face with the feeder rather than meeting it at a corner. Turf wins on a
lawn; stone wins in a cellar wall, because down there the stone is what is exposed to the shaft.

Right-click a feeder with any full block to pin that block instead. Sneak and right-click it
empty-handed to hand it back to the survey.

\begin{ietip}
It can only wear full, opaque, ordinary cubes. A feeder is a cube and always will be, so wearing a
fence would draw a fence with a solid hitbox and wearing a chest would draw a chest whose lid never
opens --- both read as a bug rather than as a disguise.
\end{ietip}

A run passes through along \textbf{one axis}, which an Engineer's Hammer turns. A feeder laid across
the grain of a run stops it exactly as a stone block would, and stops it visibly: the conduit does
not draw its arm into one it cannot enter.

\textbf{Crossing a feeder is the one place a run may change surface without a box.} The conduit on
the far side still has to have the crossing in its own plane --- the same thing a junction box asks
--- so a run can come down a north wall and continue down an east one, but it cannot continue into
conduit lying flat on the floor of the shaft.

\begin{mbplan}{A run leaving a building, seen from the side}{C = conduit on a wall, G = ground
feeder, J = junction box. The feeder fills the floor; the run crosses it and carries on below.}
\mblayer{}{J;C;G;C;J}
\end{mbplan}

A feeder \textbf{splits nothing} --- it carries the whole bundle. A breakout is a thing you should be
able to see, and a patched box wearing coloured plates is exactly that; hiding one would leave a base
whose wiring could only be found by digging.

\begin{ietip}
A feeder is scenery, not hardware. It is never a node in the wire graph, holds no energy and costs
nothing per tick --- a run crossing four of them is still one connection between the same two boxes,
four blocks longer.
\end{ietip}

\section{Redstone channels}

A face set to \textbf{read} puts the signal arriving at it onto its conductor; a face set to
\textbf{emit} puts its conductor's signal back out. The signal reaches every box on the run, so a
lever at one end of a corridor throws a lamp at the other along the same bundle already carrying the
lighting.

A face does one of the three and never two. Reading and emitting on one face is how a redstone
network latches itself on and never lets go.

\section{Capacity and cost}

Each conductor carries what a conductor carries. A bundle is not a shared pipe with an allowance to
divide up: sixteen circuits down one corridor get sixteen circuits' worth of throughput.

Sixteen conductors are \emph{one} connection in the wire graph rather than sixteen, and a junction
box carrying nothing costs one integer comparison per tick. Replacing sixteen catenaries with one
conduit costs nothing but the crafting.

Energy moves as a bucket brigade: each box hands half the difference to whichever neighbour holds
less, and drains in full into any connector on the matching breakout. The honest cost of that is one
tick per hop, and boxes exist only at corners and ends.

\chapter{Pole Signage}
\label{ch:signage}

A hundred poles across a map are a hundred identical poles. The \textbf{Utility Pole Sign} is the
tag that tells them apart: which station feeds this one, which feeder it is on, who last climbed it.

\begin{iefork}
Not in stock Immersive Engineering, and not invented either. All thirteen kinds are signs that hang
on real poles --- the shapes, the colours and what each one means are the \emph{Los Angeles
Department of Water and Power}'s and \emph{Southern California Edison}'s. They are told apart the
way the real ones are: by shape and colour, long before you are close enough to read the number.
\end{iefork}

One item, one recipe, thirteen plates. Which plate a sign is lives on the sign, not on the item, so
there is nothing to hunt for in a creative tab and nothing to craft twice.

\section{Hanging one}

Place it against the \textbf{side} of whatever holds the wires up --- a pole, a crossarm, a wall.
The plate bolts flat to that face and reads from the outside. There is no top or bottom placement: a
tag lying on a floor is not a thing that exists, and a sign whose support is broken comes down with
it, the way a torch does.

Signs have no collision box. Walking into a pole covered in them is walking into the pole.

\section{Choosing the plate}

\textbf{Hit a sign with an Engineer's Hammer} and it steps to the next of the thirteen. Keep hitting
and it walks the whole set and comes back round --- one gesture, no menu, and it works from a ladder.

\textbf{Sneak and hit it with the hammer} and the plate opens for writing:

\begin{itemize}
\item The \textbf{arrows} pick the kind, the same set the hammer walks.
\item One \textbf{text box} per line that kind carries --- one, two or three, or none at all for the
      line crossing diamond, which is a symbol rather than a label.
\item The \textbf{preview} shows the real plate with the real lettering on it, laid out by exactly
      the arithmetic the world uses. What is in that window is what ends up on the pole.
\end{itemize}

Enter and Tab step between the lines. Done and Escape both keep what was typed.

\begin{ietip}
Text is \textbf{printed to fill the plate} rather than typeset at a fixed size. A short number comes
out big and a long one comes out small, which is what the real ones do and why a pole number and a
station name can share a plate the size of a hand. The vertical strips read downwards, because a
strip six pixels wide holds a pole number no other way.
\end{ietip}

A sign taken down keeps its plate and its text, so hammering thirteen times to find the one you
meant and then mining it by accident costs nothing.

\section{What the kinds mean}

\begin{description}
\item[Parallel Generation] The horizontal red strip. Several sources are feeding the distribution
      station this pole hangs off, or power is running in a loop across the area. Mainly a LADWP
      sign, and the only one in the set that carries white text. Two lines.
\item[Pole Number] The yellow vertical strip: general pole identification, and the one every
      utility uses. One line, read downwards.
\item[Street Light] The white vertical strip. SCE hangs these for street lighting, and so do
      privately owned poles. One line, read downwards.
\item[Old Pattern] The bare metal vertical strip --- mostly replaced by something more legible, and
      still on plenty of poles. One line, read downwards, in grey.
\item[Series Light Oval] The painted white oval a series-wired street light wears: the series number
      on top and the pole number underneath.
\item[Feeder Tag] The horizontal yellow strip. For the LADWP this is the distribution station, the
      feeder off it, and how many conduit or transformer bank connections are on that connection;
      on a light pole it is the capacitor bank or the phase cutoff. Everybody else uses it for pole
      numbers.
\item[Fibre Run] The orange strip, horizontal or vertical, marking fibre and cable runs. Almost
      every utility uses these.
\item[Pole Inspection Tag] The round bolt-on tag: who inspected the pole on top, the year
      underneath. It also does duty as the tag saying where a wooden pole was logged.
\item[Tower Number] The plain yellow diamond the LADWP numbers transmission towers with, roughly
      every fourth tower. No border --- the number is painted straight onto it.
\item[Line Crossing] The same diamond with a black outline and a black cross, marking where a line
      crosses another, or a span nobody wants to walk: a valley, a ravine. It carries no text.
\item[Tower Identifier] The tower run, in two forms. The vertical one reads downwards in three
      parts --- the initials of the plant the run starts at, the tower's number, then, under a rule
      printed on the plate, the initials of the station it ends at. The horizontal one is one line,
      and is what the DC towers of the Pacific Intertie carry.
\end{description}

\section{What a sign costs to look at}

Nothing you will notice. The plate is an ordinary block model baked into the chunk mesh like any
other --- one of fifty-two, thirteen plates times four facings --- so a pole line of them costs what
a pole line of blocks costs.

Only the \emph{lettering} is drawn per frame, and only within forty-eight blocks. Past that the
plate is a pixel or two across and the number was never legible anyway, so nothing is drawn at all.
A blank sign draws nothing at any distance.

\chapter{The Hydraulic Crawler}
\label{ch:crawler}

A tracked excavator you climb into and drive, with three attachments and a demolition tool on the end
of the arm.

\begin{iefork}
Entirely this fork, and its first vehicle.
\end{iefork}

\begin{ietip}
This is \emph{not} the Excavator in chapter~\ref{ch:multiblocks}. That one is a bucket-wheel
multiblock that mines ore from a mineral vein and does not move. This one is a machine you sit in.
\end{ietip}

\section{Getting one}

Craft the Crawler and place it on level ground with about five blocks clear --- it is a large machine
and it will refuse to go anywhere it does not fit, saying so rather than silently doing nothing.
Right-click to climb in. Sneak and strike it with an Engineer's Hammer to take it away again.

\section{Driving}

\begin{itemize}
\item \textbf{W} and \textbf{S} drive; \textbf{A} and \textbf{D} steer the tracks, which turn on the
      spot the way tracks do.
\item \textbf{Looking around slews the house} --- the cab, the counterweight and the arm all turn
      with your head. You cannot look around independently of the machine, exactly as you cannot in
      the real thing.
\item \textbf{R} and \textbf{F} raise and lower the arm.
\item \textbf{X} and \textbf{Z} reach out and draw back in.
\end{itemize}

The arm is \emph{aimed} rather than jointed: one pair of keys says which way it points, the other how
far along that line the bucket sits, and the boom and stick work out their own angles. Between them
they cover a band from about two blocks out to nearly six, and from five blocks up to nearly three
below the tracks.

\begin{ietip}
The panel in the corner shows the arm's height, its reach, the fuel, and the list of controls.
\end{ietip}

\section{Attachments}

\textbf{G} changes the tool on the end of the arm, \textbf{V} works it.

\begin{itemize}
\item The \textbf{Bucket} digs and keeps what it digs, in nine slots inside the machine. Sneak and
      press \textbf{V} to tip it out; a pipe can also unload it.
\item The \textbf{Grapple} picks up one block --- or one animal --- and puts it down again. A block
      goes back exactly as it came up, facing and all.
\item The \textbf{Breaker} destroys whatever it touches and leaves it on the ground.
\end{itemize}

\begin{iewarning}
The Breaker is a demolition tool and behaves like one. It respects land claims and protected areas
exactly as your own pickaxe does --- every block it takes is on your account --- but within your own
build it will happily take a wall down faster than you can change your mind.
\end{iewarning}

\section{Fuel}

The Crawler runs on \textbf{diesel} and holds eight buckets. Refuel it by hand with any container, or
drive it up to a Gas Station Pump.

Idling costs a little and working costs a lot, so a tank goes a long way if you are only driving.
When the fuel runs low the \emph{attachment} stops but the \emph{tracks keep turning} --- that last
reserve is deliberately kept back so you can always drive home rather than being stranded next to a
half-demolished building.

\chapter{City mode}
\label{ch:citymode}

City mode trades Immersive Engineering's simulation detail for server tick time. It is aimed at city
and roleplay packs where the mod's machinery is set dressing rather than an engineering puzzle: you
keep the entire look of the build and give up the physics behind it.

\begin{iefork}
Entirely this fork. \textbf{It is off by default}, and off means byte-identical to stock.
\end{iefork}

It covers eight subsystems, each with its own switch. A subsystem is simplified only when the master
switch is on \emph{and} that subsystem has not been individually turned back off, so switching the
master off is always sufficient to restore stock behaviour.

\begin{description}
\item[Wires] One lossless push per connector instead of loss, distance weighting, a proportional
      split and a network-wide broadcast.
\item[Fluid pipes] A pipe hands its fluid to the endpoints in order rather than simulating a fill
      against each and splitting in proportion. Which of several tanks fills first is the only
      observable difference.
\item[Floodlights] Beams re-trace only when the light switches or a neighbour changes, never on a
      timer, and the light blocks one lamp may place are capped.
\item[Generators] Fuel becomes cosmetic --- a presence check and a token sip.
\item[Machines] Idle multiblocks stop re-scanning the recipe list every tick.
\item[Virtual grid] Segments switch from accounting to presence.
\item[Conduits] Bundles switch from accounting to presence.
\item[Tanks] A tank holding any fluid at all becomes an unlimited source of it. Empty is still
      empty, and filling is untouched.
\end{description}

\section{What you give up}

The numbers stop being conservation. A segment can deliver more than its feeds collected, because
nothing is being collected --- the feed is only being checked. The console's Stats tab says so
directly rather than letting the graphs be misread.

\begin{iewarning}
City mode is a decision about what your server is for. If you want the arithmetic to be honest,
leave the individual subsystem off while the master switch is on; they are independent.
\end{iewarning}

\section{The performance history}

Roughly half of everything the mod did on a profiled server turned out to be a single per-tick
network broadcast feeding wire-shock damage. That is fixed for everyone, in both modes. The full
account, with measurements, is in \texttt{docs/CITY\_MODE\_AND\_PERF.md}.


% ---------------------------------------------------------------------------
% Part IV -- reference, all of it generated
% ---------------------------------------------------------------------------
\part{Reference}

\chapter{Every block}
\label{ch:blocks}
Generated from the block metadata enums in the source tree. Metadata order is save-file order:
a block's number here is the number in the world.
\input{generated/blocks}

\chapter{Every recipe}
\label{ch:recipes}
Generated from the recipe JSON. Entries beginning with a lower-case word rather than an item name
are ore-dictionary entries, which accept any mod's equivalent.
\input{generated/recipes}

\chapter{Every configuration option}
\label{ch:config}
Generated from \texttt{Config.java}, including the comment each option ships with.
\input{generated/config}

\chapter{The in-game manual, in full}
\label{ch:ingame}
The Engineer's Manual as the book itself presents it, transcribed so the two cannot disagree.
\input{generated/ingame}

\end{document}
